closed immersion.

Finally, @@M0@@ is a torsor for f.p.p.f.; a fortiori, by f.p.q.c. descent from a ``splitting'' of the torsor, @@M1@@ is representable by a faithfully flat affine morphism of finite presentation (see: Oort, ``Commutative group schemes'', Proposition (17.4)).

\sgaStableTarget{sga6.e12.cor.3-3}\textbf{Corollary 3.3.} \uline{Under the conditions of} 3.1\uline{, we have}
\begin{enumerate}[label=\roman*)]
\item \uline{If} @@M2@@ \uline{is an isomorphism, then} @@M3@@ \uline{is an isomorphism.}
\item \uline{Supposing that the} @@M4@@\uline{-scheme} @@M5@@ \uline{is affine, if} @@M6@@ \uline{is surjective, then} @@M7@@ \uline{is surjective.}
\end{enumerate}

Indeed, (i) follows from the following diagram of sheaves for the Zariski topology
@@M8@@
by applying the five lemma to it.

(ii) follows from the following diagram with exact rows
@@M9@@
because the outer vertical arrows are isomorphisms, @@M10@@ being affine.

